\subsection{Propiedades funtoriales de los invariantes diferenciales}
\label{IV.16.4}

\begin{env}[16.4.1]
\label{IV.16.4.1}
Consideremos un diagrama conmutativo de morfismos de preesquemas
\[
  \xymatrix{
    X \ar[d]_-{f} & X' \ar[l]_-u \ar[d]^-{f'} \\
    S & S' \ar[l]^-w
  }
  \tag{16.4.1.1}
  \label{IV.16.4.1.1}
\]
De él se deduce un diagrama conmutativo
\[
  \xymatrix{
    X \ar[d]_-{\Delta_f} & X' \ar[l]_-u \ar[d]^-{\Delta_{f'}} \\
    X\times_S X & X'\times_{S'}X' \ar[l]^-v
  }
\]
donde $v$ es el morfismo compuesto (I, 5.3.5 y 5.3.15)
\[
  X'\times_{S'}X'
  \xrightarrow{(p'_1,p'_2)_S}X'\times_S X'
  \xrightarrow{u\times_S u}X\times_S X.
  \tag{16.4.1.2}
  \label{IV.16.4.1.2}
\]
Se deducen, por tanto, de $u$ y $v$, como se explicó en (16.2.1),
homomorfismos de haces de anillos aumentados
\[
  \nu_n:\rho^*(\mathscr{P}_{X/S}^n)\longrightarrow
  \mathscr{P}_{X'/S'}^n
  \tag{16.4.1.3}
  \label{IV.16.4.1.3}
\]
(donde se ha puesto $u=(\rho,\lambda)$); estos homomorfismos forman un sistema
proyectivo y dan, por tanto, en el límite un homomorfismo de haces de anillos
aumentados
\[
  \nu_\infty:\rho^*(\mathscr{P}_{X/S}^\infty)\longrightarrow
  \mathscr{P}_{X'/S'}^\infty;
  \tag{16.4.1.4}
  \label{IV.16.4.1.4}
\]
por otra parte, al pasar a los cocientes, los homomorfismos $\nu_n$ dan un
dihomomorfismo de álgebras graduadas (relativo a $\lambda^\#$):
\[
  \operatorname{gr}(u):\rho^*(\mathscr{G}r_\bullet(\mathscr{P}_{X/S}))
  \longrightarrow\mathscr{G}r_\bullet(\mathscr{P}_{X'/S'}).
  \tag{16.4.1.5}
  \label{IV.16.4.1.5}
\]
\end{env}

\begin{env}[16.4.2]
\label{IV.16.4.2}
\oldpage[IV-4]{17}
Si se tiene un diagrama conmutativo
\[
  \xymatrix{
    X \ar[d]_-{f} & X' \ar[l]_-u \ar[d]^-{f'} &
      X'' \ar[l]_-{u'} \ar[d]^-{f''} \\
    S & S' \ar[l]^-w & S'' \ar[l]^-{w'}
  }
\]
se deduce un diagrama conmutativo
\[
  \xymatrix{
    X \ar[d]_-{\Delta_f} & X' \ar[l]_-u \ar[d]^-{\Delta_{f'}} &
      X'' \ar[l]_-{u'} \ar[d]^-{\Delta_{f''}} \\
    X\times_S X & X'\times_{S'}X' \ar[l]^-v &
      X''\times_{S''}X'' \ar[l]^-{v'}
  }
\]
donde $v'$ se define a partir de $u'$, $w'$, $f'$, $f''$ como $v$ a partir de
$u$, $w$, $f$, $f'$. Se verifica inmediatamente que, si $u''=u\circ u'$,
$w''=w\circ w'$, entonces el morfismo compuesto $v\circ v'$ es igual al morfismo
$v''$ definido a partir de $u''$, $w''$, $f$, $f''$ como $v$ a partir de $u$,
$w$, $f$, $f'$. Si se pone $u'=(\rho',\lambda')$,
$u''=(\rho'',\lambda'')$, resulta entonces de (16.2.1) que el homomorfismo
$\nu_n'':\rho''^*(\mathscr{P}_{X/S}^n)\to\mathscr{P}_{X''/S''}^n$ es igual al
compuesto
\[
  \rho'^*(\rho^*(\mathscr{P}_{X/S}^n))
  \xrightarrow{\rho'^*(\nu_n^\#)}
  \rho'^*(\mathscr{P}_{X'/S'}^n)
  \xrightarrow{\nu_n'}\mathscr{P}_{X''/S''}^n
\]
y se tienen propiedades de transitividad análogas para los homomorfismos
(16.4.1.4) y (16.4.1.5), lo que permite decir que los
$\mathscr{P}_{X/S}^n$, $\mathscr{P}_{X/S}^\infty$ y
$\mathscr{G}r_\bullet(\mathscr{P}_{X/S})$ dependen funtorialmente de $f$.
\end{env}

\begin{env}[16.4.3]
\label{IV.16.4.3}
Se verifica inmediatamente (por ejemplo, reduciéndose al caso afín con ayuda de
(16.3.7)) que, con las notaciones de (16.4.1), el diagrama
\[
  \xymatrix{
    \rho^*(\mathcal{O}_X) \ar[r]^{\lambda^\#} \ar[d] &
      \mathcal{O}_{X'} \ar[d] \\
    \rho^*(\mathscr{P}_{X/S}^n) \ar[r]_{\nu_n} &
      \mathscr{P}_{X'/S'}^n
  }
  \tag{16.4.3.1}
  \label{IV.16.4.3.1}
\]
donde las flechas verticales son las que definen las estructuras de álgebra
elegidas en (16.3.5) (es decir, las que provienen de las primeras proyecciones),
es conmutativo; lo mismo ocurre con el diagrama
\[
  \xymatrix{
    \rho^*(\mathcal{O}_X) \ar[r]^{\lambda^\#}
      \ar[d]_{\rho^*(d_{X/S}^n)} &
      \mathcal{O}_{X'} \ar[d]^{d_{X'/S'}^n} \\
    \rho^*(\mathscr{P}_{X/S}^n) \ar[r]_{\nu_n} &
      \mathscr{P}_{X'/S'}^n
  }
  \tag{16.4.3.2}
  \label{IV.16.4.3.2}
\]
\oldpage[IV-4]{18}
las flechas verticales definen aquí las estructuras de álgebra provenientes de las
segundas proyecciones; además, si $\sigma$ y $\sigma'$ son las simetrías
canónicas correspondientes a $f$ y $f'$ (16.3.4), se tiene
\[
  \nu_n\circ\rho^*(\sigma)=\sigma'\circ\nu_n
\]
que permite pasar de uno de los diagramas precedentes al otro. Se deduce, por tanto, de
(16.4.3.1) un homomorfismo canónico de $\mathcal{O}_{X'}$-álgebras aumentadas
\[
  P^n(u):u^*(\mathscr{P}_{X/S}^n)
  =\mathscr{P}_{X/S}^n\otimes_{\mathcal{O}_X}\mathcal{O}_{X'}
  \longrightarrow\mathscr{P}_{X'/S'}^n
  \tag{16.4.3.3}
  \label{IV.16.4.3.3}
\]
y resulta de (16.4.3.2) que el diagrama
\[
  \xymatrix{
    \mathcal{O}_{X'} \ar[r]^{\operatorname{id}}
      \ar[d]_{u^*(d_{X/S}^n)} &
      \mathcal{O}_{X'} \ar[d]^{d_{X'/S'}^n} \\
    u^*(\mathscr{P}_{X/S}^n) \ar[r]_{P^n(u)} &
      \mathscr{P}_{X'/S'}^n
  }
  \tag{16.4.3.4}
  \label{IV.16.4.3.4}
\]
es conmutativo. Se deduce un homomorfismo de $\mathcal{O}_{X'}$-álgebras
graduadas
\[
  \operatorname{Gr}_\bullet(u):u^*(\mathscr{G}r_\bullet(\mathscr{P}_{X/S}))
  \longrightarrow\mathscr{G}r_\bullet(\mathscr{P}_{X'/S'})
  \tag{16.4.3.5}
  \label{IV.16.4.3.5}
\]
y, en particular, un homomorfismo de $\mathcal{O}_{X'}$-módulos
\[
  \operatorname{Gr}_1(u):
  \Omega_{X/S}^1\otimes_{\mathcal{O}_X}\mathcal{O}_{X'}
  \longrightarrow\Omega_{X'/S'}^1
  \tag{16.4.3.6}
  \label{IV.16.4.3.6}
\]
que da lugar a un diagrama conmutativo
\[
  \xymatrix{
    \mathcal{O}_{X'} \ar[r]^{\operatorname{id}}
      \ar[d]_{d_{X/S}\otimes 1} &
      \mathcal{O}_{X'} \ar[d]^{d_{X'/S'}} \\
    \Omega_{X/S}^1\otimes_{\mathcal{O}_X}\mathcal{O}_{X'} \ar[r] &
      \Omega_{X'/S'}^1
  }
  \tag{16.4.3.7}
  \label{IV.16.4.3.7}
\]
\end{env}

\begin{env}[16.4.4]
\label{IV.16.4.4}
Cuando $S=\operatorname{Spec}(A)$, $S'=\operatorname{Spec}(A')$,
$X=\operatorname{Spec}(B)$, $X'=\operatorname{Spec}(B')$ son afines, de modo
que se tiene un diagrama conmutativo de homomorfismos de anillos
\[
  \xymatrix{
    B \ar[r] & B' \\
    A \ar[u] \ar[r] & A' \ar[u]
  }
\]
la imagen de $\mathfrak{I}_{B/A}$ en $B'\otimes_{A'}B'$ está contenida en
$\mathfrak{I}_{B'/A'}$, y el homomorfismo $\nu_n$ corresponde al
homomorfismo de anillos $P_{B/A}^n\to P_{B'/A'}^n$ deducido del homomorfismo
$B\otimes_A B\to B'\otimes_{A'}B'$ al pasar a los cocientes. El homomorfismo
(16.4.3.6) corresponde al homomorfismo definido en (0, 20.5.4.1), y el
diagrama conmutativo (16.4.3.7) al diagrama (0, 20.5.4.2).
\end{env}

\oldpage[IV-4]{19}
\begin{proposition}[16.4.5]
\label{IV.16.4.5}
Supongamos que $X'=X\times_S S'$, siendo $f'$ y $u$ las proyecciones canónicas.
Entonces los homomorfismos canónicos $P^n(u)$ (16.4.3.3) y
$\operatorname{Gr}_1(u)$ (16.4.3.6) son biyectivos.
\end{proposition}

En efecto, se tiene entonces $X'\times_{S'}X'=(X\times_S X)\times_S S'$, y se puede, por tanto,
aplicar (16.2.3, (ii)) reemplazando $g$ por la primera proyección
$p_1:X\times_S X\to X$ y $f$ por la diagonal $\Delta_f$.

Obsérvese también que, bajo las hipótesis de (16.4.5), el homomorfismo
$\operatorname{Gr}_\bullet(u)$ (16.4.3.5) es sobreyectivo, pero no biyectivo en
general. Sin embargo (16.2.4):

\begin{corollary}[16.4.6]
\label{IV.16.4.6}
Siendo las hipótesis las de (16.4.5), supongamos además que $w:S'\to S$ es
plano (resp. que los $\mathscr{G}r_n(\mathscr{P}_{X/S})$ son
$\mathcal{O}_X$-módulos planos para $n\leq m$); entonces el homomorfismo
\[
  \operatorname{Gr}_n(u):u^*(\mathscr{G}r_n(\mathscr{P}_{X/S}))
  \longrightarrow\mathscr{G}r_n(\mathscr{P}_{X'/S'})
\]
es biyectivo para todo $n$ (resp. para $n\leq m$).
\end{corollary}

En efecto, si $w$ es plano, lo mismo ocurre con
$v:X'\times_{S'}X'\to X\times_S X$, por lo que la conclusión resulta de (16.2.4).

\begin{env}[16.4.7]
\label{IV.16.4.7}
Sean $S$ un preesquema, $\mathscr{E}$ un $\mathcal{O}_S$-módulo
cuasicoherente, y pongamos $X=\mathbf{V}(\mathscr{E})$ (II, 1.7.8), fibrado
vectorial asociado a $\mathscr{E}$, igual a
$\operatorname{Spec}(\mathscr{S}_{\mathcal{O}_S}(\mathscr{E}))$. Sea
$f:X\to S$ el morfismo estructural. Para todo abierto $U$ de $S$, y toda
sección $t\in\Gamma(U,\mathscr{E})$, $t$ se identifica con una sección de
$\mathscr{S}_{\mathcal{O}_S}(\mathscr{E})$ sobre $U$; sea $t'$ su
imagen en
$\Gamma(f^{-1}(U),\mathcal{O}_X)=\Gamma(U,f_*(\mathcal{O}_X))=
\Gamma(U,\mathscr{S}_{\mathcal{O}_S}(\mathscr{E}))$, y pongamos
\[
  \delta(t)=d_{X/S}^n(t')-t'\in
  \Gamma(f^{-1}(U),\mathscr{P}_{X/S}^n);
  \tag{16.4.7.1}
  \label{IV.16.4.7.1}
\]
está claro que $\delta$ es un dihomomorfismo de módulos (correspondiente al
homomorfismo de anillos
$\Gamma(U,\mathcal{O}_S)\to\Gamma(f^{-1}(U),\mathcal{O}_X)$) de
$\Gamma(U,\mathscr{E})$ en $\Gamma(f^{-1}(U),\mathscr{P}_{X/S}^n)$, cuya
imagen pertenece además al ideal de aumentación de
$\Gamma(f^{-1}(U),\mathscr{P}_{X/S}^n)$. Se deduce (haciendo variar $U$)
un homomorfismo canónico de $\mathcal{O}_X$-álgebras
\[
  f^*(\mathscr{S}_{\mathcal{O}_S}(\mathscr{E}))
  \longrightarrow\mathscr{P}_{X/S}^n
  \tag{16.4.7.2}
  \label{IV.16.4.7.2}
\]
y, según la observación precedente, si $\mathscr{K}$ es el ideal que es núcleo de
la aumentación $\mathscr{S}_{\mathcal{O}_S}(\mathscr{E})\to\mathcal{O}_S$,
la imagen de $\mathscr{K}^{n+1}$ por (16.4.7.2) es nula, de modo que, al
factorizar por $\mathscr{K}^{n+1}$, se obtiene finalmente un homomorfismo canónico
\[
  \delta_n:f^*(\mathscr{S}_{\mathcal{O}_S}(\mathscr{E})/
  \mathscr{K}^{n+1})\longrightarrow\mathscr{P}_{X/S}^n.
  \tag{16.4.7.3}
  \label{IV.16.4.7.3}
\]
\end{env}

\begin{proposition}[16.4.8]
\label{IV.16.4.8}
Bajo las condiciones de (16.4.7), los homomorfismos $\delta_n$ son biyectivos y
forman un sistema proyectivo de isomorfismos; se deduce un isomorfismo de
$\mathcal{O}_X$-álgebras graduadas
\[
  f^*(\mathscr{S}_{\mathcal{O}_S}^{\bullet}(\mathscr{E}))
  \longrightarrow\mathscr{G}r_\bullet(\mathscr{P}_{X/S}).
  \tag{16.4.8.1}
  \label{IV.16.4.8.1}
\]
\end{proposition}

El hecho de que los homomorfismos (16.4.7.3) formen un sistema proyectivo resulta
inmediatamente de su definición. Para probar que son isomorfismos, bastará
\oldpage[IV-4]{20}
probar que (16.4.8.1) es un isomorfismo, pues las filtraciones de los dos miembros de
(16.4.7.3) son finitas (Bourbaki, \emph{Alg. comm.}, cap. III, \S~2,
n\textsuperscript{o} 8, cor. 3 del teo. 1). Para ello, consideremos la sucesión exacta
escindida de $\mathcal{O}_S$-módulos
\[
  \xymatrix{
    0 \ar[r] & \mathscr{E} \ar[r]^{u} &
    \mathscr{E}\oplus\mathscr{E} \ar[r]^{v} &
    \mathscr{E} \ar[r] & 0
  }
  \tag{16.4.8.2}
  \label{IV.16.4.8.2}
\]
donde, para todo par de secciones $s$, $t$ de $\mathscr{E}$ sobre un abierto
$U$ de $S$, se toma $u(s)=(-s,s)$ y $v(s,t)=s+t$. Se tiene
\[
  X\times_S X=\operatorname{Spec}
  (\mathscr{S}_{\mathcal{O}_S}(\mathscr{E})\otimes_{\mathcal{O}_S}
  \mathscr{S}_{\mathcal{O}_S}(\mathscr{E}))
  =\operatorname{Spec}(\mathscr{S}_{\mathcal{O}_S}
  (\mathscr{E}\oplus\mathscr{E}))
\]
(II, 1.4.6 y 1.7.11), y el morfismo diagonal $X\to X\times_S X$ corresponde
(II, 1.2.7) al homomorfismo de $\mathcal{O}_S$-álgebras
$\mathscr{S}(v):\mathscr{S}_{\mathcal{O}_S}(\mathscr{E}\oplus\mathscr{E})
\to\mathscr{S}_{\mathcal{O}_S}(\mathscr{E})$ (II, 1.7.4), de modo que, si
$\mathscr{I}$ es el núcleo de este homomorfismo, se tiene
\[
  \mathscr{P}_{X/S}^n=f^*(\mathscr{S}_{\mathcal{O}_S}
  (\mathscr{E}\oplus\mathscr{E})/\mathscr{I}^{n+1}).
\]

La proposición será consecuencia del lema siguiente:

\begin{lemma}[16.4.8.3]
\label{IV.16.4.8.3}
Sean $Y$ un espacio anillado,
$0\to\mathscr{F}'\xrightarrow{u}\mathscr{F}\xrightarrow{v}\mathscr{F}''\to0$
una sucesión exacta de $\mathcal{O}_Y$-módulos tal que todo punto $y\in Y$
tiene una vecindad abierta $V$ tal que la sucesión
$0\to\mathscr{F}'|V\to\mathscr{F}|V\to\mathscr{F}''|V\to0$ es escindida.
Sea $\mathscr{I}$ el ideal que es núcleo de
$\mathscr{S}(v):\mathscr{S}_{\mathcal{O}_Y}(\mathscr{F})\to
\mathscr{S}_{\mathcal{O}_Y}(\mathscr{F}'')$, y sea
$\operatorname{gr}_{\mathscr{I}}^\bullet
(\mathscr{S}_{\mathcal{O}_Y}(\mathscr{F}))$ la $\mathcal{O}_Y$-álgebra graduada
asociada a la $\mathcal{O}_Y$-álgebra
$\mathscr{S}_{\mathcal{O}_Y}(\mathscr{F})$ dotada de la filtración
$\mathscr{I}$-preádica. Entonces el homomorfismo de $\mathcal{O}_Y$-álgebras
graduadas
\[
  \mathscr{S}_{\mathcal{O}_Y}^{\bullet}(\mathscr{F}')
  \otimes_{\mathcal{O}_Y}
  \mathscr{S}_{\mathcal{O}_Y}^{\bullet}(\mathscr{F}'')
  \longrightarrow
  \operatorname{gr}_{\mathscr{I}}^\bullet
  (\mathscr{S}_{\mathcal{O}_Y}(\mathscr{F}))
  \tag{16.4.8.4}
  \label{IV.16.4.8.4}
\]
(donde el primer miembro es el producto tensorial graduado de las
$\mathcal{O}_Y$-álgebras simétricas dotadas de su graduación canónica
(II, 1.7.4 y 2.1.2)), proveniente de la inyección canónica
\[
  \mathscr{F}'\longrightarrow\mathscr{I}
  =\operatorname{gr}_{\mathscr{I}}^1
  (\mathscr{S}_{\mathcal{O}_Y}(\mathscr{F})),
\]
es biyectivo.
\end{lemma}

En efecto, la inyección $\mathscr{F}'\to\mathscr{I}$ da canónicamente un
homomorfismo de $\mathcal{O}_Y$-álgebras graduadas\linebreak
$\mathscr{S}_{\mathcal{O}_Y}^{\bullet}(\mathscr{F}')\to\allowbreak
\operatorname{gr}_{\mathscr{I}}^\bullet
(\mathscr{S}_{\mathcal{O}_Y}(\mathscr{F}))$, y como el segundo miembro es por
definición una $\mathscr{S}_{\mathcal{O}_Y}^{\bullet}(\mathscr{F}'')$-álgebra
graduada, se deduce el homomorfismo canónico (16.4.8.4) tensorizando el
precedente con $\mathscr{S}_{\mathcal{O}_Y}^{\bullet}(\mathscr{F}'')$. Para
probar el lema, como la cuestión es local, basta considerar el caso en que
$\mathscr{F}=\mathscr{F}'\oplus\mathscr{F}''$, siendo $u$ y $v$ los
homomorfismos canónicos. Entonces el álgebra graduada
$\mathscr{S}_{\mathcal{O}_Y}^{\bullet}(\mathscr{F})$ se identifica canónicamente con el
producto tensorial graduado
$\mathscr{S}_{\mathcal{O}_Y}^{\bullet}(\mathscr{F}')
\otimes_{\mathcal{O}_Y}
\mathscr{S}_{\mathcal{O}_Y}^{\bullet}(\mathscr{F}'')$ (II, 1.7.4), y es
inmediato que $\mathscr{I}$ es entonces el ideal
$\mathscr{I}'\otimes_{\mathcal{O}_Y}
\mathscr{S}_{\mathcal{O}_Y}^{\bullet}(\mathscr{F}'')$, donde $\mathscr{I}'$ es
el ideal de aumentación de
$\mathscr{S}_{\mathcal{O}_Y}^{\bullet}(\mathscr{F}')$, es decir, la suma
(directa) de los $\mathscr{S}_{\mathcal{O}_Y}^m(\mathscr{F}')$ para $m\geq1$. Se
concluye que $\mathscr{I}^n=\mathscr{I}'^n\otimes_{\mathcal{O}_Y}
\mathscr{S}_{\mathcal{O}_Y}^{\bullet}(\mathscr{F}'')$, donde esta vez
$\mathscr{I}'^n$ es la suma (directa) de los
$\mathscr{S}_{\mathcal{O}_Y}^m(\mathscr{F}')$ para $m\geq n$; se tiene por consiguiente
$\mathscr{I}^n/\mathscr{I}^{n+1}=
\mathscr{S}_{\mathcal{O}_Y}^n(\mathscr{F}')\otimes_{\mathcal{O}_Y}
\mathscr{S}_{\mathcal{O}_Y}^{\bullet}(\mathscr{F}'')$, lo que prueba que
(16.4.8.4) es biyectivo.

\oldpage[IV-4]{21}
Demostrado este lema, queda por ver que el homomorfismo (16.4.8.1) es
efectivamente la imagen por $f^*$ del homomorfismo (16.4.8.4) correspondiente a la sucesión
exacta (16.4.8.2); se comprueba fácilmente que esto resulta de la definición de
$u$ (16.4.8.2) y de la de $\delta$ (16.4.7.1), teniendo en cuenta la definición
de la estructura de $\mathcal{O}_X$-álgebra sobre $\mathscr{P}_{X/S}^n$ y de
la de $d_{X/S}^n$ (16.3.5 y 16.3.6).

En particular:

\begin{corollary}[16.4.9]
\label{IV.16.4.9}
Bajo las condiciones de (16.4.7), se tiene un isomorfismo canónico
\[
  \operatorname{gr}_1(\delta):f^*(\mathscr{E})
  \xrightarrow{\sim}\Omega_{X/S}^1.
  \tag{16.4.9.1}
  \label{IV.16.4.9.1}
\]
\end{corollary}

\begin{corollary}[16.4.10]
\label{IV.16.4.10}
Si $S=\operatorname{Spec}(A)$, $\mathscr{E}=\mathcal{O}_S^m$, de modo que
\[
  X=\operatorname{Spec}(A[T_1,\ldots,T_m]),
\]
$\mathscr{P}_{X/S}^n$ se identifica canónicamente con el
$\mathcal{O}_X$-álgebra correspondiente al $A[T_1,\ldots,T_m]$-álgebra cociente
\[
  A[T_1,\ldots,T_m,U_1,\ldots,U_m]/\mathfrak{K}^{n+1},
\]
donde las $U_i$ ($1\leq i\leq m$) son $m$ nuevas indeterminadas y
$\mathfrak{K}$ es el ideal engendrado por $U_1,\ldots,U_m$.
\end{corollary}

Se recupera así, en particular, la estructura de $\Omega_{X/S}^1$ en este caso
(0, 20.5.13).

Obsérvese además que $d_{X/S}^n$ hace corresponder entonces a un polinomio
$F(T_1,\ldots,T_m)$ la clase mód.\ $\mathfrak{K}^{n+1}$ de
$F(T_1+U_1,\ldots,T_m+U_m)$, como resulta de la definición (16.4.7.1).

\begin{proposition}[16.4.11]
\label{IV.16.4.11}
Sean $f:X\to S$ un morfismo, $g:S\to X$ una sección sobre $S$ de $X$,
$S_g^{(n)}$ la $n$-ésima vecindad infinitesimal de $S$ para la inmersión $g$
(16.1.2). Existe entonces un único isomorfismo de
$\mathcal{O}_S$-álgebras
\[
  \varpi_n:g^*(\mathscr{P}_{X/S}^n)\longrightarrow
  \mathcal{O}_{S_g^{(n)}}
  \tag{16.4.11.1}
  \label{IV.16.4.11.1}
\]
(para la estructura de $\mathcal{O}_S$-álgebra sobre
$\mathcal{O}_{S_g^{(n)}}$ definida por $f$ (16.1.7)), que hace conmutativo el
diagrama
\[
  \xymatrix{
    \mathcal{O}_S=g^*(\mathcal{O}_X)
      \ar[rr]^-{\lambda_n}\ar[dr]_-{g^*(d_{X/S}^n)}
      &&\mathcal{O}_{S_g^{(n)}}\\
    &g^*(\mathscr{P}_{X/S}^n)\ar[ur]_-{\varpi_n}
  }
  \tag{16.4.11.2}
  \label{IV.16.4.11.2}
\]
(donde $\lambda_n$ es el homomorfismo estructural).
\end{proposition}

En virtud de (I, 5.3.7), donde se reemplazan $X$, $Y$, $S$ por $S$, $X$, $S$,
respectivamente, y $f$ por $g$, los diagramas
\[
  \xymatrix@=3pc{
    S\ar[r]^-g\ar[d]_-g&X\ar[d]^-{\Delta_f}
      &S\ar[r]^-g\ar[d]_-g&X\ar[d]^-{\Delta_f}\\
    X\ar[r]_{(g\circ f,1_X)_S}&X\times_S X
      &X\ar[r]_{(1_X,g\circ f)_S}&X\times_S X
  }
  \tag{16.4.11.3}
  \label{IV.16.4.11.3}
\]
\oldpage[IV-4]{22}
identifican $S$ con el producto fibrado de los $(X\times_S X)$-preesquemas $X$ y $X$ para los
morfismos $\Delta_f$ y $(g\circ f,1_X)_S$ (resp.\ $(1_X,g\circ f)_S$).
Por otra parte, los diagramas
\[
  \xymatrix@=3pc{
    X\ar[r]^{(g\circ f,1_X)_S}\ar[d]_-f&X\times_S X\ar[d]^-{p_1}
      &X\ar[r]^{(1_X,g\circ f)_S}\ar[d]_-f&X\times_S X\ar[d]^-{p_2}\\
    S\ar[r]_{g}&X&S\ar[r]_{g}&X
  }
  \tag{16.4.11.4}
  \label{IV.16.4.11.4}
\]
identifican $X$ con el producto fibrado de los $X$-preesquemas $S$ y $X\times_S X$ para los
morfismos $g$ y $p_1$ (resp.\ $p_2$) (caso particular de la fórmula
de asociatividad (I, 3.3.9.1)). Puede decirse que $\Delta_f$, considerado como
sección sobre $X$ de $X\times_S X$ (relativamente a $p_1$ o $p_2$), desempeña el papel
de una \emph{sección universal} para las secciones sobre $S$ de $X$: cada una de estas
secciones $g$ se deduce, en efecto, de ella por el \emph{cambio de base}
$(g\circ f,1_X)_S:X\to X\times_S X$. La definición del homomorfismo
$\varpi_n$ y el hecho de que sea biyectivo resultan, por tanto, de estas observaciones y de
(16.2.3, (ii)) aplicado al primer diagrama (16.4.11.4). La conmutatividad del
diagrama (16.4.11.2) resulta asimismo de (16.2.3, (ii)) aplicado esta vez al
segundo diagrama (16.4.11.4). Para explicitar $\varpi_n$, basta considerar el
caso en que $g$ es una inmersión \emph{cerrada}: en efecto, para todo $s\in S$, hay
una vecindad abierta $W$ de $s$ en $S$ tal que $g(W)$ es cerrado en un
abierto $U$ de $X$, y está claro que $g|W$ es una sección sobre $W$ del morfismo
$U\cap f^{-1}(W)\to W$, restricción de $f$, y $g(W)$ es \emph{a fortiori}
cerrado en $U\cap f^{-1}(W)$. Se puede, por tanto, suponer que $S$ es un
subpreesquema cerrado de $X$ definido por un ideal cuasicoherente $\mathscr{K}$.
Entonces las definiciones precedentes muestran que, si $W$ es un abierto de $S$,
$t$ una sección de $\mathcal{O}_X$ sobre $f^{-1}(W)$,
$\varpi_n(d^n t|W)$ es igual a la imagen canónica de $t$ en
$\Gamma(W,(\mathcal{O}_X/\mathscr{K}^{n+1})|W)$. La unicidad de $\varpi_n$
resulta, por tanto, de que la imagen de $\mathcal{O}_X$ por $d_{X/S}^n$ genera el
$\mathcal{O}_X$-módulo $\mathscr{P}_{X/S}^n$ (16.3.8).

\begin{corollary}[16.4.12]
\label{IV.16.4.12}
Sean $k$ un cuerpo, $X$ un $k$-preesquema, $x$ un punto de $X$ racional sobre
$k$. Entonces $(\mathscr{P}_{X/k}^n)_x\otimes_{\mathcal{O}_x}k(x)$ es
canónicamente isomorfo (como $k(x)$-álgebra aumentada) a
$\mathcal{O}_x/\mathfrak{m}_x^{n+1}$.
\end{corollary}

Basta considerar la única sección sobre $k$, $g$, de $X$ tal que
$g(\operatorname{Spec}(k))=\{x\}$.

\begin{corollary}[16.4.13]
\label{IV.16.4.13}
Sean $f:X\to S$ un morfismo, $s$ un punto de $S$,
$X_s=X\times_S\operatorname{Spec}(k(s))$ la fibra de $f$ en $s$. Si
$x\in X_s$ es racional sobre $k(s)$,
$(\mathscr{P}_{X/S}^n)_x\otimes_{\mathcal{O}_s}k(s)$ es canónicamente
isomorfo a $\mathcal{O}_{X_s,x}/\mathfrak{m}'_x{}^{\,n+1}$, donde
$\mathfrak{m}'_x$ es el ideal maximal de $\mathcal{O}_{X_s,x}$; más
precisamente, este isomorfismo hace corresponder a $(d^n t)_x\otimes1$ (donde $t$ es
una sección de $\mathcal{O}_X$ sobre una vecindad abierta de $x$ en
$X$) la clase de $t_x\otimes1$ módulo $\mathfrak{m}'_x{}^{\,n+1}$.
\end{corollary}

Esto resulta de (16.4.5) y (16.4.12).

Los corolarios precedentes justifican la terminología de «haz de partes
principales de orden $n$».

\begin{proposition}[16.4.14]
\label{IV.16.4.14}
Sea $\rho:A\to B$ un homomorfismo de anillos, $S$ un subconjunto
multiplicativamente estable de $B$. Entonces los homomorfismos canónicos
\[
  S^{-1}P_{B/A}^n\longrightarrow P_{S^{-1}B/A}^n
  \tag{16.4.14.1}
  \label{IV.16.4.14.1}
\]
\oldpage[IV-4]{23}
deducidos de los homomorfismos canónicos
$P_{B/A}^n\to P_{S^{-1}B/A}^n$ (16.4.4), forman un sistema proyectivo y son
biyectivos.
\end{proposition}

Basta observar que
\[
  S^{-1}((B\otimes_A B)/\mathfrak{I}^{n+1})
  =S^{-1}(B\otimes_A B)/(S^{-1}\mathfrak{I})^{n+1}
\]
por planitud, y que
$S^{-1}(B\otimes_A B)=(S^{-1}B)\otimes_A(S^{-1}B)$ (I, 1.3.4).

\begin{corollary}[16.4.15]
\label{IV.16.4.15}
Con las notaciones de (16.4.14), sea $R$ un subconjunto multiplicativo de
$A$ tal que $\rho(R)\subset S$. Entonces se tienen isomorfismos canónicos
\[
  S^{-1}P_{B/A}^n\xrightarrow{\sim}P_{S^{-1}B/R^{-1}A}^n
  \tag{16.4.15.1}
  \label{IV.16.4.15.1}
\]
que forman un sistema proyectivo.
\end{corollary}

Basta, evidentemente, definir isomorfismos canónicos
\[
  P_{S^{-1}B/A}^n\xrightarrow{\sim}P_{S^{-1}B/R^{-1}A}^n,
  \tag{16.4.15.2}
  \label{IV.16.4.15.2}
\]
es decir, se reduce al caso en que $\rho(R)$ está formado por elementos
\emph{invertibles} de $B$. Pero entonces el isomorfismo (16.4.15.2) se
deduce simplemente del isomorfismo canónico
$B\otimes_A B\to B\otimes_{R^{-1}A}B$ al pasar a los cocientes
(0$_{\rm I}$, 1.5.3).

\begin{corollary}[16.4.16]
\label{IV.16.4.16}
Sean $f:X\to S$ un morfismo de preesquemas, $x$ un punto de $X$,
$s=f(x)$. Entonces se tienen isomorfismos canónicos
\[
  (\mathscr{P}_{X/S}^n)_x\xrightarrow{\sim}
  P_{\mathcal{O}_x/\mathcal{O}_s}^n
  \tag{16.4.16.1}
  \label{IV.16.4.16.1}
\]
que forman un sistema proyectivo.
\end{corollary}

Se deducen de aquí isomorfismos para los graduados asociados y, en particular,
un isomorfismo canónico
\[
  (\Omega_{X/S}^1)_x\xrightarrow{\sim}
  \Omega_{\mathcal{O}_x/\mathcal{O}_s}^1.
  \tag{16.4.16.2}
  \label{IV.16.4.16.2}
\]

\begin{corollary}[16.4.17]
\label{IV.16.4.17}
Sean $k$ un cuerpo, $K$ el cuerpo de fracciones racionales
$k(T_1,\ldots,T_r)$. Entonces, para todo entero $n$, el homomorfismo de
$K[U_1,\ldots,U_r]$ ($U_i$ indeterminadas) en $P_{K/k}^n$ que hace corresponder a cada $U_i$
$d^nT_i-T_i\cdot1$ es sobreyectivo y define un isomorfismo
del cociente $K[U_1,\ldots,U_r]/\mathfrak{m}^{n+1}$ (donde $\mathfrak{m}$ es
el ideal engendrado por las $U_i$) sobre $P_{K/k}^n$.
\end{corollary}

Esto resulta de (16.4.8), (16.4.10) y (16.4.14), donde se toma $A=k$,
$B=k[T_1,\ldots,T_r]$ y $S=B-\{0\}$.

Se recupera así el hecho de que los $dT_i$ forman una base del $K$-espacio
vectorial $\Omega_{K/k}^1$ (0, 20.5.10).

\begin{proposition}[16.4.18]
\label{IV.16.4.18}
Sean $f:X\to Y$, $g:Y\to Z$ dos morfismos de preesquemas; consideremos
los homomorfismos canónicos de $\mathcal{O}_X$-álgebras aumentadas
(16.4.3.3)
\[
  g_{X/Y/Z}:\mathscr{P}_{X/Z}^n\longrightarrow\mathscr{P}_{X/Y}^n
  \tag{16.4.18.1}
  \label{IV.16.4.18.1}
\]
\[
  f_{X/Y/Z}:f^*(\mathscr{P}_{Y/Z}^n)
  \longrightarrow\mathscr{P}_{X/Z}^n.
  \tag{16.4.18.2}
  \label{IV.16.4.18.2}
\]
Entonces $g_{X/Y/Z}$ es sobreyectivo, y su núcleo es el ideal engendrado por la imagen
por $f_{X/Y/Z}$ del ideal de aumentación de
$f^*(\mathscr{P}_{Y/Z}^n)$.
\end{proposition}

\oldpage[IV-4]{24}
Observemos primero que $g_{X/Y/Z}$ corresponde al caso en que, en (16.4.3.3), se toma
$X'=X$, $S'=Y$, $S=Z$ y $u=1_X$, y $f_{X/Y/Z}$ al caso en que se reemplazan
$X'$, $X$, $S$, $S'$ por $X$, $Y$, $Z$, $Z$, respectivamente, y $u$, $f$ por
$f$, $g$, respectivamente.

Se tiene un diagrama conmutativo (I, 5.3.5)
\[
  \xymatrix{
    X\ar[r]^-{\Delta_f}\ar[dr]_-f
      &X\times_Y X\ar[d]^-p\ar[r]^-j
      &X\times_Z X\ar[d]^-{f\times_Z f}\\
    &Y\ar[r]_-{\Delta_g}&Y\times_Z Y
  }
  \tag{16.4.18.3}
  \label{IV.16.4.18.3}
\]
donde $j=(1_X,1_X)_Z$ es una inmersión, $j\circ\Delta_f=\Delta_{g\circ f}$ y
$p$ es el morfismo estructural. Como basta considerar el caso en que $X$, $Y$,
$Z$ son afines, se puede suponer que las inmersiones $\Delta_f$, $\Delta_g$ y
$j$ son cerradas, de modo que $\mathcal{O}_X$ y
$\mathcal{O}_{X\times_Y X}$ se identifican, respectivamente, con
$\mathcal{O}_{X\times_Z X}/\mathscr{J}$ y
$\mathcal{O}_{X\times_Z X}/\mathscr{L}$, donde
$\mathscr{J}\supset\mathscr{L}$ son los dos ideales cuasicoherentes
correspondientes, respectivamente, a las inmersiones $\Delta_{g\circ f}$ y $j$. El
$\mathcal{O}_X$-álgebra $\mathscr{P}_{X/Z}^n$ se identifica, por tanto, con
$\mathcal{O}_{X\times_Z X}/\mathscr{J}^{n+1}$, y
$\mathscr{P}_{X/Y}^n$ se identifica con
$\mathcal{O}_{X\times_Y X}/(\mathscr{J}/\mathscr{L})^{n+1}$, es decir, con
$\mathcal{O}_{X\times_Z X}/(\mathscr{J}^{n+1}+\mathscr{L})$, y por consiguiente con el
cociente de $\mathscr{P}_{X/Z}^n$ por
$(\mathscr{J}^{n+1}+\mathscr{L})/\mathscr{J}^{n+1}$. Pero se sabe
(\emph{loc. cit.}) que $p$ y $j$ hacen de $X\times_Y X$ el producto fibrado de los
$(Y\times_Z Y)$-preesquemas $Y$ y $X\times_Z X$; por tanto, si $\mathcal{O}_Y$ se
identifica con $\mathcal{O}_{Y\times_Z Y}/\mathscr{K}$, donde $\mathscr{K}$ es
el ideal correspondiente a $\Delta_g$, $\mathscr{L}$ es igual a
$(f\times_Z f)^*(\mathscr{K})\cdot\mathcal{O}_{X\times_Z X}$ (I, 4.4.5).
Como $(\mathscr{J}^{n+1}+\mathscr{L})/\mathscr{J}^{n+1}$ es el ideal de
$\mathscr{P}_{X/Z}^n$ engendrado por la imagen de $\mathscr{L}$, se deduce la
proposición.

\begin{corollary}[16.4.19]
\label{IV.16.4.19}
Con las notaciones de (16.4.18), se tiene una sucesión exacta de
$\mathcal{O}_X$-módulos cuasicoherentes
\[
  \xymatrix{
    f^*(\Omega_{Y/Z}^1)\ar[r]^-{f_{X/Y/Z}}&
    \Omega_{X/Z}^1\ar[r]^-{g_{X/Y/Z}}&
    \Omega_{X/Y}^1\ar[r]&0.
  }
  \tag{16.4.19.1}
  \label{IV.16.4.19.1}
\]
\end{corollary}

Cuando $X$, $Y$, $Z$ son afines, se recupera así la sucesión exacta
(0, 20.5.7.1).

\begin{proposition}[16.4.20]
\label{IV.16.4.20}
Sean $f:Y\to Z$ un morfismo, $j:X\to Y$ una inmersión cerrada,
$\mathscr{K}$ el ideal cuasicoherente de $\mathcal{O}_Y$ correspondiente a $j$.
Entonces se tiene
$\mathscr{P}_{X/Y}^n=\mathcal{O}_X=\mathcal{O}_Y/\mathscr{K}$,
el homomorfismo canónico
$j_{X/Y/Z}:j^*(\mathscr{P}_{Y/Z}^n)\to\mathscr{P}_{X/Z}^n$ es sobreyectivo, y
su núcleo es el ideal de $j^*(\mathscr{P}_{Y/Z}^n)$ engendrado por
$j^*(\mathcal{O}_Y\cdot d_{Y/Z}^n(\mathscr{K}))$ (debe observarse que
$d_{Y/Z}^n(\mathscr{K})$ es un subhaz de grupos conmutativos de
$\mathscr{P}_{Y/Z}^n$, pero no un $\mathcal{O}_Y$-módulo en general).
\end{proposition}

Se sabe (I, 5.3.8) que la diagonal
$\Delta_j:X\to X\times_Y X$ es un isomorfismo, de donde se sigue la primera afirmación.
Si $\varpi_1$ y $\varpi_2$ son los dos homomorfismos de álgebras
$\mathcal{O}_Y\to\mathscr{P}_{Y/Z}^n$ correspondientes, respectivamente, a las dos
proyecciones canónicas $p_1$, $p_2$ de $Y\times_Z Y$ sobre $Y$, recordemos que,
por definición (16.3.5 y 16.3.6), $\varpi_1$ es el homomorfismo estructural del
$\mathcal{O}_Y$-álgebra $\mathscr{P}_{Y/Z}^n$ y
$\varpi_2=d_{Y/Z}^n$. El $\mathcal{O}_X$-álgebra
$j^*(\mathscr{P}_{Y/Z}^n)$ se identifica, por tanto, con
$\mathscr{P}_{Y/Z}^n/\varpi_1(\mathscr{K})\mathscr{P}_{Y/Z}^n$ y su cociente
por el ideal engendrado por $j^*(d_{Y/Z}^n(\mathscr{K}))$ con
\[
  \mathscr{P}_{Y/Z}^n/
  (\varpi_1(\mathscr{K})+\varpi_2(\mathscr{K}))\mathscr{P}_{Y/Z}^n.
\]
Observemos ahora que se tiene un diagrama conmutativo
\oldpage[IV-4]{25}
\[
  \xymatrix{
    Y\ar[d]_-{\Delta_f}&X\ar[l]_-j\ar[d]^-{\Delta_{f\circ j}}\\
    Y\times_Z Y&X\times_Z X\ar[l]^-{j\times_Z j}
  }
\]
que identifica $X$ con el producto fibrado de los $(Y\times_Z Y)$-preesquemas $Y$ y
$X\times_Z X$ (I, 5.3.7). Como $j\times_Z j$ es una inmersión, se deduce
de esta observación y de (16.2.2) que, si $\Delta_{Y/Z}^n$ y
$\Delta_{X/Z}^n$ denotan las vecindades infinitesimales de orden $n$ de $Y$
y $X$ para las inmersiones canónicas $\Delta_f$ y $\Delta_{f\circ j}$,
respectivamente, se tiene un diagrama
\[
  \xymatrix{
    \Delta_{Y/Z}^n\ar[d]&\Delta_{X/Z}^n\ar[l]\ar[d]\\
    Y\times_Z Y&X\times_Z X\ar[l]^-{j\times_Z j}
  }
\]
que hace de $\Delta_{X/Z}^n$ el producto fibrado de los $(Y\times_Z Y)$-preesquemas
$\Delta_{Y/Z}^n$ y $X\times_Z X$. También puede decirse que
$\mathscr{P}_{X/Z}^n$ se identifica con el haz de anillos
$\mathscr{P}_{Y/Z}^n\otimes_{\mathcal{O}_{Y\times_Z Y}}
\mathcal{O}_{X\times_Z X}$. Pero se ve inmediatamente (por ejemplo, reduciéndose
al caso afín) que
\[
  \mathcal{O}_{X\times_Z X}=
  \mathcal{O}_{Y\times_Z Y}/
  (p_1^*(\mathscr{K})+p_2^*(\mathscr{K}))
  \mathcal{O}_{Y\times_Z Y}.
\]
Por tanto, $\mathscr{P}_{X/Z}^n$ se identifica con el cociente de
$\mathscr{P}_{Y/Z}^n$ por el ideal engendrado por la imagen en
$\mathscr{P}_{Y/Z}^n$ de $p_1^*(\mathscr{K})+p_2^*(\mathscr{K})$. Pero, por
definición, este ideal también está engendrado por
$\varpi_1(\mathscr{K})+\varpi_2(\mathscr{K})$. Q.E.D.

\begin{corollary}[16.4.21]
\label{IV.16.4.21}
Sean $f:Y\to Z$ un morfismo, $j:X\to Y$ una inmersión. Se tiene una sucesión
exacta de $\mathcal{O}_X$-módulos cuasicoherentes
\[
  \xymatrix{
    \mathscr{N}_{X/Y}\ar[r]&j^*(\Omega_{Y/Z}^1)\ar[r]&
    \Omega_{X/Z}^1\ar[r]&0.
  }
  \tag{16.4.21.1}
  \label{IV.16.4.21.1}
\]
\end{corollary}

Cuando $X$, $Y$ y $Z$ son afines, se recupera así la sucesión exacta
(0, 20.5.12.1).

\begin{corollary}[16.4.22]
\label{IV.16.4.22}
Si $f:X\to S$ es un morfismo localmente de presentación finita,
$\mathscr{P}_{X/S}^n$ y $\Omega_{X/S}^1$ son
$\mathcal{O}_X$-módulos cuasicoherentes de presentación finita.
\end{corollary}

Se reduce inmediatamente al caso en que $S=\operatorname{Spec}(A)$ es afín,
$X=\operatorname{Spec}(B)$, donde
$B=A[T_1,\ldots,T_r]/\mathfrak{K}$, siendo $\mathfrak{K}$ un ideal de tipo
finito de $C=A[T_1,\ldots,T_r]$. Se aplica entonces (16.4.20) con $Z=S$,
$Y=\operatorname{Spec}(C)$ y $\mathscr{K}=\widetilde{\mathfrak{K}}$. Entonces
$j^*(\mathscr{P}_{Y/Z}^n)$ es un $\mathcal{O}_X$-módulo libre de rango finito
(16.4.10), y la hipótesis sobre $\mathfrak{K}$ implica que
$j^*(\mathcal{O}_Y\cdot d_{Y/Z}^n(\mathscr{K}))$ genera un
$\mathcal{O}_X$-módulo cuasicoherente de tipo finito; de donde se sigue la conclusión.

\begin{proposition}[16.4.23]
\label{IV.16.4.23}
Sean $X$, $Y$ dos $S$-preesquemas, $Z=X\times_S Y$ su producto fibrado,
$p:X\times_S Y\to X$ y $q:X\times_S Y\to Y$ las proyecciones canónicas.
Entonces el homomorfismo canónico
\[
  p_{Z/X/S}\oplus q_{Z/Y/S}:
  p^*(\Omega_{X/S}^1)\oplus q^*(\Omega_{Y/S}^1)
  \longrightarrow\Omega_{(X\times_S Y)/S}^1
  \tag{16.4.23.1}
  \label{IV.16.4.23.1}
\]
es biyectivo.
\end{proposition}

\oldpage[IV-4]{26}
El diagrama conmutativo
\[
  \xymatrix{
    Y\ar[d]_-g&X\times_S Y\ar[l]_-q\ar[d]_-h&
      X\times_S Y\ar[l]_-{\operatorname{id}}\ar[d]^-p\\
    S&S\ar[l]^-{\operatorname{id}}&X\ar[l]^-f
  }
\]
da una factorización del \emph{isomorfismo} canónico $P^n(p)$ (16.4.5)
\[
  p^*(\mathscr{P}_{X/S}^n)\longrightarrow
  \mathscr{P}_{Z/S}^n\longrightarrow\mathscr{P}_{Z/Y}^n
\]
y, del mismo modo, intercambiando los papeles de $X$ y $Y$, se tiene una factorización
del \emph{isomorfismo} $P^n(q)$
\[
  q^*(\mathscr{P}_{Y/S}^n)\longrightarrow
  \mathscr{P}_{Z/S}^n\longrightarrow\mathscr{P}_{Z/X}^n.
\]
Esto prueba que el homomorfismo canónico (16.4.18.1)
\[
  p_{Z/X/S}:p^*(\mathscr{P}_{X/S}^n)\longrightarrow\mathscr{P}_{Z/S}^n
  \quad\text{(resp.\
  $q_{Z/Y/S}:q^*(\mathscr{P}_{Y/S}^n)\to\mathscr{P}_{Z/S}^n$)}
\]
es \emph{inyectivo}, y que el núcleo del homomorfismo canónico sobreyectivo
(16.4.18.2)
\[
  \mathscr{P}_{Z/S}^n\longrightarrow\mathscr{P}_{Z/Y}^n
  \quad\text{(resp.\
  $\mathscr{P}_{Z/S}^n\to\mathscr{P}_{Z/X}^n$)}
\]
es suplementario de la imagen de $p_{Z/X/S}$ (resp.\ $q_{Z/Y/S}$). Pero,
por otra parte, este núcleo está, en virtud de (16.4.18), engendrado por la imagen por
$q_{Z/Y/S}$ (resp.\ $p_{Z/X/S}$) del ideal de aumentación de
$q^*(\mathscr{P}_{Y/S}^n)$ (resp.\ $p^*(\mathscr{P}_{X/S}^n)$). Se concluye
la proposición considerando el caso $n=1$.

Se generaliza inmediatamente (16.4.23) al caso de un producto fibrado de un número finito
arbitrario de $S$-preesquemas.

\begin{remarks}[16.4.24]
\label{IV.16.4.24}
\begin{enumerate}
\item[(i)] Veremos (17.2.3) que, cuando el morfismo $f:X\to Y$ de
(16.4.18) es \emph{liso}, el homomorfismo $f_{X/Y/Z}$ de (16.4.19.1) es
localmente \emph{invertible por la izquierda} y, en particular, inyectivo. Del mismo modo,
cuando el morfismo $f\circ j:X\to Z$ de (16.4.20) es \emph{liso},
el homomorfismo de la izquierda en (16.4.21.2) es localmente
\emph{invertible por la izquierda} y \emph{a fortiori} inyectivo (17.2.5). En el cap. V,
daremos también una variante, para módulos sobre
preesquemas, de los «módulos de imperfección» estudiados en (0, 20.6), y de las sucesiones
exactas en las que figuran.

\item[(ii)] Sean $X$ un espacio topológico, $\mathscr{A}$ un haz
de anillos sobre $X$ y $\mathscr{B}$ una $\mathscr{A}$-álgebra sobre $X$. Entonces está
claro que
\[
  U\longmapsto
  P_{\Gamma(U,\mathscr{B})/\Gamma(U,\mathscr{A})}^n
  \qquad (U\text{ abierto en }X)
\]
es un prehaz de $\Gamma(U,\mathscr{B})$-álgebras aumentadas; por tanto, el
haz asociado $\mathscr{P}_{\mathscr{B}/\mathscr{A}}^n$ es una
$\mathscr{B}$-álgebra aumentada. En el caso particular en que $X$ es un
preesquema, $f=(\psi,\theta):X\to S$ un morfismo de preesquemas, resulta
fácilmente de (16.4.16) y de la exactitud del funtor $\varinjlim$ que
$\mathscr{P}_{X/S}^n$ es canónicamente isomorfo a
$\mathscr{P}_{\mathcal{O}_X/\psi^*(\mathcal{O}_S)}^n$. De ello se sigue que el
formalismo desarrollado en el presente parágrafo podría considerarse como
un
\oldpage[IV-4]{27}
caso particular de un formalismo diferencial para espacios anillados dotados
de un haz de álgebras sobre el haz estructural. Sin embargo, no hemos
querido partir de este punto de vista, menos intuitivo y menos cómodo para las
aplicaciones. Parece, además, que, para las diversas especies de
«variedades», la construcción «global» de los $\mathscr{P}^n$, análoga a la que
utilizamos aquí, está también mejor adaptada a las aplicaciones.
\end{enumerate}
\end{remarks}
